\section{Method}

\subsection{Problem Formulation}

We consider hallucination in large language models as a decision-making problem under uncertainty, rather than purely a generation problem.

Given an input query $x$ and contextual information $C = \{c_1, c_2, \dots, c_k\}$, a language model generates a response $y$. 
The objective is to produce responses that are both grounded and useful, while avoiding unsupported or misleading outputs.

\subsubsection{Hallucination as Decision-Making}

We define hallucination as a failure to make appropriate decisions under uncertainty. 
Specifically, the model must jointly determine: 
(i) whether to answer or abstain, 
(ii) how to construct the response, and 
(iii) how confident the response should be.

Under this formulation, hallucination occurs when the model produces responses that are either unsupported by available information or expressed with excessive confidence under uncertainty.

\subsubsection{Controlled Contextual Scenarios}

To systematically study model behavior, we consider two controlled scenarios:

\begin{itemize}
    \item \textbf{No reliable information:} The context does not contain sufficient evidence to answer the query. The desired behavior is to abstain.
    \item \textbf{Partially correct information:} The context contains a mixture of relevant and irrelevant information, including at least one correct supporting evidence. The desired behavior is to generate a selective and grounded response.
\end{itemize}

These scenarios enable explicit modeling of the decision boundary between answering and abstaining.

\subsubsection{Groundedness and Uncertainty}

We introduce a groundedness score $s \in [0,1]$, measuring the degree to which a response is supported by the available context, following the groundedness metric proposed in RAG evaluation frameworks \cite{es2024ragas}.

We quantify model uncertainty using the average token-level entropy, which serves as a proxy for predictive uncertainty and has been widely adopted in prior work~\cite{farquhar2024detecting}:

\[
H_{\text{avg}} = -\frac{1}{T} \sum_{t=1}^{T} \sum_{v \in \mathcal{V}} p_t(v)\log p_t(v)
\]

where $T$ is the response length and $p_t(v)$ denotes the probability of token $v$ at step $t$.

We interpret low entropy as high confidence and high entropy as high uncertainty. 
Hallucination often arises when the model exhibits low entropy despite low groundedness.

\subsubsection{Objective}

Our objective is to learn a policy that:
(i) abstains when reliable information is unavailable, 
(ii) produces grounded responses when sufficient evidence exists, 
(iii) avoids overconfident behavior under uncertainty, and 
(iv) maintains response quality and usefulness.

This formulation naturally leads to a reinforcement learning framework that optimizes behavior under uncertainty rather than purely maximizing likelihood.


\subsection{Reward Function}

We design a behavior-oriented reward function that explicitly models hallucination as a decision-making problem under uncertainty. 
The reward integrates four components: behavior alignment, confidence control, response quality, and length regularization.

\[
R = R_{\text{behavior}} + R_{\text{confidence}} + R_{\text{quality}} + R_{\text{length}}
\]

\subsubsection{Behavior Alignment}

We distinguish two controlled scenarios: (i) no reliable information, where the model should abstain, and (ii) partially correct information, where the model should answer.

Let $a \in \{0,1\}$ denote the model action, where $a=1$ indicates abstention and $a=0$ indicates answering.

\[
R_{\text{behavior}} =
\begin{cases}
+R_{\text{abstain}}, & \text{if abstention is correct} \\
-\alpha, & \text{if abstention is incorrect} \\
0, & \text{otherwise}
\end{cases}
\]

\subsubsection{Confidence Control via Entropy}

We introduce an overconfidence penalty based on entropy:

\[
H_{\text{avg}} = -\frac{1}{T} \sum_{t=1}^{T} \sum_{v \in \mathcal{V}} p_t(v)\log p_t(v)
\]

\[
\text{overconf} = \max(0, U_{\min} - H_{\text{avg}})
\]

\[
R_{\text{confidence}} = -\gamma (1 - s)\cdot \text{overconf}
\]

This discourages overly confident predictions when groundedness is low, while preserving confident generation when responses are well-supported.

\subsubsection{Response Quality Shaping}

We introduce a quality score $q \in [0,1]$, computed from multiple dimensions including completeness, actionability, evidence use, and conciseness, following LLM-based evaluation approaches such as G-Eval~\cite{geval}.

\[
R_{\text{quality}} = \beta \cdot q
\]

This component promotes informative and well-structured responses beyond factual correctness.

\subsubsection{Length Regularization}

To prevent degenerate short responses, we impose a length constraint:

\[
L \in [L_{\text{low}}, \infty]
\]

\[
R_{\text{length}} =
-\eta \left(
\max\left(0, \frac{L_{\text{low}} - L}{L_{\text{low}}}\right)
\right)
\]

This penalizes under-informative outputs. In practice, we set $L_{\text{low}} = 50$.

\subsubsection{Summary}

The reward function jointly optimizes decision-making, uncertainty-aware behavior, response quality, and generation stability, enabling balanced behavior under uncertainty.


\subsection{Training Framework}

We adopt a two-stage LoRA training framework consisting of preference-based initialization followed by reinforcement learning refinement.

\subsubsection{Stage 1: Preference Optimization (DPO)}

We first perform Direct Preference Optimization (DPO)~\cite{dpo} to obtain a stable initialization.

Given $(x, C)$, we construct preference pairs $(y^{+}, y^{-})$:

\begin{itemize}
    \item $y^{+}$: preferred response (grounded or correct)
    \item $y^{-}$: less preferred response (hallucinated or incorrect)
\end{itemize}

\[
\mathcal{L}_{\text{DPO}} = - \log \sigma \left( \beta_{\text{dpo}} \left( \log \pi_{\theta}(y^{+}|x,C) - \log \pi_{\theta}(y^{-}|x,C) \right) \right)
\]

\subsubsection{Stage 2: Reinforcement Learning (GRPO)}

We refine the model using Group Relative Policy Optimization (GRPO)~\cite{GRPO}.

For each $(x, C)$, we sample:

\[
\{y_1, y_2, \dots, y_N\} \sim \pi_{\theta}(\cdot | x, C)
\]

Each response is evaluated using the reward $R(y)$.

\[
A_i = R(y_i) - \frac{1}{N} \sum_{j=1}^{N} R(y_j)
\]

The policy is updated to increase the likelihood of higher-reward responses relative to the group.

\subsubsection{Training Stability}

We observe that reinforcement learning may lead to degenerate strategies such as overly short responses. 
To mitigate this, we incorporate quality-based reward shaping and length regularization to maintain a balance between hallucination mitigation and response usefulness.